% Generated draft: requires mathematical and source review.
\documentclass[11pt]{article}
\usepackage[margin=25mm]{geometry}
\usepackage{fontspec}
\setmainfont{Latin Modern Roman}
\usepackage{amsmath,amssymb,mathtools,unicode-math}
\setmathfont{Latin Modern Math}
\usepackage{enumitem,needspace,xurl,hyperref}
\setlength{\parindent}{0pt}
\setlength{\parskip}{5pt}
\setlength{\emergencystretch}{3em}
\newcommand{\N}{\mathbb{N}}
\newcommand{\Z}{\mathbb{Z}}
\newcommand{\Q}{\mathbb{Q}}
\newcommand{\R}{\mathbb{R}}
\newcommand{\C}{\mathbb{C}}
\newcommand{\F}{\mathbb{F}}
\newcommand{\A}{\mathbb{A}}
\newcommand{\PP}{\mathbb{P}}
\newcommand{\E}{\mathbb{E}}
\newcommand{\eps}{\varepsilon}
\newcommand{\dd}{\,\mathrm{d}}
\begin{document}
\section*{UnsolvedMath ID: 30006024}
\textbf{Source problem code:} catalogue entry
\subsection{Definitions and mathematical statement}
A hyperbolic surface is a surface equipped with a Riemannian metric of constant curvature \(-1\). The supplied record does not specify orientability, compactness, completeness, boundary components, or cusps. For a compact connected closed surface \(X\), a \emph{one-face map} means a finite connected graph cellularly embedded in \(X\), with exactly one face: its complement is homeomorphic to an open disk. This definition states the usual closed-surface convention; the record does not specify whether a different convention for surfaces with boundary or cusps is intended.

The term \emph{spine} is not defined in the supplied record. Its geometric construction, any auxiliary choices, and whether it retains edge lengths or other decorations must be specified before its law can be determined. In particular, interpreting a spine as a graph onto which the entire surface deformation retracts would not apply to a compact closed surface.

The source poses the following questions.
\begin{enumerate}
\item Determine the asymptotic law, and in particular any limiting law, of the spine of a random hyperbolic surface of genus \(g\) as \(g\to\infty\). To make this precise, one must specify a probability measure \(\mu_g\) on a specified space of genus-\(g\) hyperbolic surfaces, sample \(X_g\sim\mu_g\), define its spine \(S(X_g)\), and specify a measurable encoding and any normalization \(T_g\) into a common topological state space \(\mathcal E\). A possible precise limiting-law target would then be to identify a probability measure \(\nu\) such that
\[
\mathcal L\bigl(T_g(S(X_g))\bigr)\Longrightarrow\nu
\qquad (g\to\infty),
\]
where \(\mathcal L\) denotes probability law and \(\Longrightarrow\) denotes weak convergence on the chosen state space. None of \(\mu_g\), \(S\), \(T_g\), \(\mathcal E\), or the intended topology is specified in the record, so this display is a formulation template, not a supplied conjecture asserting that a limit exists.
\item Find other natural geometric constructions of one-face maps on hyperbolic surfaces. The source does not define the admissible surfaces, the construction against which these are to be considered ``other,'' or a mathematical criterion for naturalness.
\item Find a geometric adaptation of Chapuy's bijection that constructs random surfaces. This question appears explicitly in the original statement although it is omitted from the shortened statement. The intended bijection, decorations, geometric data, and probability distribution on the resulting surfaces are not specified.
\end{enumerate}

The supplied surrounding exposition describes Chapuy's bijection between one-face maps and decorated trees, cited there as reference [4], as a key tool in the proof of reference [6]. The bibliography entries for [4] and [6] are absent from this record. Consequently, that contextual description does not determine a precise bijection or theorem here.
\subsection{Short English statement}
Determine the large-genus distribution of a geometrically defined spine of a random hyperbolic surface, find natural geometric constructions of one-face maps, and seek a geometric adaptation of Chapuy's bijection for constructing random surfaces. A precise probability model, spine construction, convergence framework, and target bijection remain unspecified. The supplied record labels the questions open; this draft does not independently verify that status.
\subsection{Sources}
\begin{itemize}
\item Supplied local UnsolvedMath registry record, permanent ID \(30006024\), title \emph{Spines and One-Face Maps of Random Hyperbolic Surfaces}, source problem code \texttt{catalogue entry}, in the Oberwolfach Reports Open Problems collection. No individual proposer is supplied; the recorded proposed year is 2024.
\item Supplied primary-source citation: \emph{Riemann Surfaces: Random, Flat, and Hyperbolic Geometry} (2025), Oberwolfach Reports, pp.\ 2369--2424, DOI \texttt{10.4171/owr/2024/41}. \url{https://doi.org/10.4171/owr/2024/41}. The original excerpt includes the page marker 2405. The supplied citation uses publication year 2025, while the DOI and proposed-year metadata use 2024; these are retained without resolving their relationship.
\item Supplied literature source labeled ``Mirzakhani--Petri'': \url{https://arxiv.org/abs/1710.09727}. The record supplies no title or specific result from this reference.
\item Literature assessment dated exactly \texttt{2026-08-22}: the supplied curation labels the questions open and states that existing bijections and random-surface results do not answer them in the stated generality. This is supplied dated curation, not an independently verified literature assessment. The background explicitly requests independent verification of time-sensitive status and cited claims.
\item The registry labels the statement \texttt{corrected\_verified} and reports that the official Oberwolfach report was checked, the mathematical question isolated, and notation repaired. These are supplied source-verification claims, not checks performed for this draft. The excerpt's references [4] and [6] lack bibliography entries in the supplied record.
\end{itemize}
\end{document}
